\textbf{Purpose:} Crowd-vote supervision can alter confidence ranking without consistently improving category choice. We examine this distinction when a facial-expression predictor omits categories in the annotation vocabulary.
\textbf{Methods:} Fifteen matched Swin-T runs compare hard, tie-aware hard, dominant-confidence-preserving uniform-minority, tie-preserving uniform-minority, and full soft targets on FER2013 images with FER+ votes. Exact-pixel-isolated selection, calibration and testing support common-coverage, fixed-image and fixed-epoch comparisons. Complete-vote disagreement is separated into unsupported mass, supported ambiguity and category-choice excess. Three adaptive phases reuse one image source; they are not independent confirmation.
\textbf{Results:} At 80\% coverage, selected-model Soft-minus-Uniform disagreement is -0.38 percentage points; Soft-minus-Tie-Uniform is -0.59. On Uniform-selected images, Soft instead differs by +0.09 points. Changing the reference on identical majority-subset images substantially changes F1. Soft retains more zero-supported-mass images than Uniform. All 36 nonempty empirical test operating points for the original four conditions exceed their calibration targets.
\textbf{Conclusion:} Target representation, accepted-image composition, categorical reference and vocabulary mismatch jointly shape apparent reliability. The observed within-source contrasts support diagnostic reporting, not a novel architecture, universal superiority, deployment risk guarantee or validated messaging benefit.